\FormSection{1. General information:}
\FormIndented{Project code: \ProjectCode}
\FormIndented{Project title: \ProjectTitleEnglish}
\FormIndented{Coordinator: \PrincipalName}
\FormIndented{Implementing institution: \PrincipalUnit}
\FormIndented{Duration: from March 2026 to August 2026.}
\FormSection{2. Objective(s):}
Develop a prototype for analyzing and consolidating flood-rescue reports under intermittent connectivity, combining edge inference, compact transmission and event clustering.
\FormSection{3. Creativeness and innovativeness:}
Combine adaptive sending, persistent queues and event clustering with conditional edge-localization and exact-copy invariance audits. Evaluate clustering, ranking and dispatch separately.
\FormSection{4. Research results:}
The seed-1024 image model achieves \NewImageAccuracy\% accuracy and \NewImageMacroFOne\% macro-F1 on 256 test images, with three critical errors. A separate split manifest and runtime/compression/mobile benchmarks for this model remain unavailable. Historical benchmarks are reported separately; PTQ/QAT summary--CSV discrepancies remain unresolved. Structured urgency uses a synthetic 32-state table rather than a Vietnamese text corpus. RQ1--RQ3 do not establish general superiority or field rescue benefits.
\FormSection{5. Products:}
Flutter application, FastAPI/SQLite backend, map dashboard, model assets, experiment data, research report and bilingual bulletins/summaries. An ISDS 2026 draft exists without publication or formal adoption evidence.
\FormSection{6. Effects, technology transfer means and applicability:}
The prototype supports education, research and workflow demonstration. Planned transfer uses the application, source code, dashboard and operating documentation; no completed transfer record is available. Text modeling, urgency-field integration, repeated Android/network/SMS tests and expert validation remain necessary before operational deployment.
